\documentclass[11pt,a4paper]{article}
\usepackage[utf8]{inputenc}
\usepackage[T1]{fontenc}
\usepackage[spanish,es-minimal,es-noshorthands]{babel}
\usepackage{lmodern}
\usepackage[sfdefault,lining]{FiraSans}
\usepackage[a4paper,top=2.3cm,bottom=2.2cm,left=1.8cm,right=1.8cm,headheight=1.3cm,headsep=0.4cm,footskip=1cm]{geometry}
\usepackage{amsmath,amssymb}
\usepackage{icomma}
\usepackage{graphicx}
\usepackage[table]{xcolor}
\usepackage{tikz}
\usepackage{fontawesome5}
\usepackage{ccicons}
\usepackage{fancyhdr}
\usepackage{lastpage}
\usepackage{tasks}
\usepackage{needspace}
\usepackage[most]{tcolorbox}
\usepackage{array}
\usepackage[hidelinks]{hyperref}
\AtBeginDocument{\hypersetup{pdftitle={Ficha de ejercicios de la unidad 5: Trigonometría. 4º ESO},pdfsubject={Matemáticas B 4º ESO},pdfkeywords={trigonometría, radianes, seno, coseno, tangente, relaciones fundamentales, triángulos rectángulos, circunferencia goniométrica, reducción al primer cuadrante}}}

\definecolor{navy}{HTML}{1B2748}
\definecolor{accent}{HTML}{E8603A}
\definecolor{teal}{HTML}{1F9E90}
\definecolor{soft}{HTML}{EEF1F8}
\definecolor{muted}{HTML}{5D6685}

\setlength{\parindent}{0pt}
\setlength{\parskip}{3pt}
\settasks{label=\alph*), label-format=\bfseries\color{navy}, label-width=1.3em, item-indent=1.8em, column-sep=0.6em, after-item-skip=5pt, before-skip=3pt, after-skip=4pt}

% ---------- capçalera i peu ----------
\newcommand{\llicencia}{\href{https://creativecommons.org/licenses/by-nc-sa/4.0/deed.es}{\ccbyncsa}\ \ Licencia Creative Commons Atribución-NoComercial-CompartirIgual 4.0 Internacional (CC BY-NC-SA 4.0)}
\pagestyle{fancy}
\fancyhf{}
\fancyhead[L]{\raisebox{-0.35\height}{\includegraphics[height=0.95cm]{logo-gva.png}}}
\fancyhead[C]{\small\color{muted}Ficha de ejercicios · Unidad 5: Trigonometría · 4\textsuperscript{o} ESO B}
\fancyhead[R]{\raisebox{-0.35\height}{\includegraphics[height=0.95cm]{logo-vermellar.png}}}
\fancyfoot[L]{\scriptsize\color{muted}\llicencia}
\fancyfoot[R]{\small\color{muted}\thepage}
\renewcommand{\headrulewidth}{0.4pt}
\renewcommand{\headrule}{\color{muted!40}\hrule height \headrulewidth}
\fancypagestyle{primera}{\fancyhf{}\renewcommand{\headrulewidth}{0pt}
  \fancyfoot[L]{\scriptsize\color{muted}\llicencia}
  \fancyfoot[R]{\small\color{muted}\thepage}}

% ---------- nivells i exercicis ----------
\newcommand{\estrelles}[1]{{\color{accent}\ifcase#1\or\faStar\faStar[regular]\faStar[regular]\or\faStar\faStar\faStar[regular]\or\faStar\faStar\faStar\fi}}
\newcounter{exer}
\newcommand{\ex}[1]{\par\needspace{4\baselineskip}\medskip\refstepcounter{exer}\noindent{\color{navy}\bfseries\theexer.}\ \ {\small\estrelles{#1}}\ \ }

\newcommand{\seccio}[2]{\par\needspace{16\baselineskip}\bigskip
  \begin{tcolorbox}[colback=navy,colframe=navy,coltext=white,arc=2mm,boxrule=0pt,left=3mm,right=3mm,top=1.2mm,bottom=1.2mm]
  {\large\bfseries #1\quad #2}
  \end{tcolorbox}\vspace{-2pt}}

\newcommand{\exf}[1]{\par\medskip\refstepcounter{exer}\noindent{\color{navy}\bfseries\theexer.}\ \ {\small\estrelles{#1}}\ \ }
\newcolumntype{C}[1]{>{\centering\arraybackslash}p{#1}}
% quadre de recordatori (només a la fitxa d'exercicis)
\newcommand{\recorda}[1]{\par\smallskip\begin{tcolorbox}[colback=accent!7,colframe=accent!55,arc=2mm,boxrule=0.5pt,left=2.5mm,right=2.5mm,top=1mm,bottom=1mm]\small{\color{accent}\faLightbulb[regular]}\ \textbf{Recuerda:}\ #1\end{tcolorbox}}
\usetikzlibrary{calc}
\usepackage{enumitem}
\providecommand{\tg}{\operatorname{tg}}
\newcommand{\res}[1]{#1}\newcommand{\rest}[1]{#1}

\def\fmneed{11}% que el títol d'una secció no quede sol al final del full
% =====================================================================
% Capçalera FentMates de les fitxes: un color per a cada curs.
% Els colors són els de la pàgina inicial (LEVELS a assets/js/materials.js);
% el logo ve de _eines/logo/pdf/ (python3 _eines/logo/genera.py).
%
% Ús, al final del preàmbul de la unitat:
%   % =====================================================================
% Capçalera FentMates de les fitxes: un color per a cada curs.
% Els colors són els de la pàgina inicial (LEVELS a assets/js/materials.js);
% el logo ve de _eines/logo/pdf/ (python3 _eines/logo/genera.py).
%
% Ús, al final del preàmbul de la unitat:
%   % =====================================================================
% Capçalera FentMates de les fitxes: un color per a cada curs.
% Els colors són els de la pàgina inicial (LEVELS a assets/js/materials.js);
% el logo ve de _eines/logo/pdf/ (python3 _eines/logo/genera.py).
%
% Ús, al final del preàmbul de la unitat:
%   \input{../comu/fentmates.tex}
%   \fmcurs{1eso}
% i, al començament del document:
%   \capcalera{Matemàtiques}{1r ESO}{Fitxa d'exercicis · Unitat 2}{Divisibilitat}
% Opcional abans de \input: \def\fmneed{14} (espai mínim abans de cada secció).
% =====================================================================
\makeatletter
\def\fm@logodir{../../logo/pdf}
\providecommand{\fmneed}{10}
\def\fm@cursnom{}

% clau del curs, color i color del text damunt del color
\newcommand{\fm@curs}[3]{\@namedef{fm@color@#1}{#2}\@namedef{fm@text@#1}{#3}}
\fm@curs{1eso}{E6B84A}{navy}
\fm@curs{2eso}{3B82F6}{white}
\fm@curs{3eso}{FF7A50}{white}
\fm@curs{4eso-a}{A78BFA}{white}
\fm@curs{4eso-b}{6D28D9}{white}
\fm@curs{1batx-mi}{9B1C1C}{white}
\fm@curs{1batx-ccss}{EF6B6B}{white}
\fm@curs{2batx-mii}{1E6B47}{white}
\fm@curs{2batx-ccss}{2BB3A3}{white}

\newcommand{\fmcurs}[1]{%
  \@ifundefined{fm@color@#1}{\PackageError{fentmates}{Curs desconegut: #1}{}}{}%
  \edef\fm@tmp{\noexpand\definecolor{curs}{HTML}{\@nameuse{fm@color@#1}}%
               \noexpand\colorlet{curstext}{\@nameuse{fm@text@#1}}}\fm@tmp
  \colorlet{cursfosc}{curs!72!black}%
  \def\fm@logo{\fm@logodir/vertical-#1.pdf}%
  \def\fm@marca{\fm@logodir/marca-#1.pdf}}

% ---------- capçalera de la primera pàgina ----------
% #1 matèria (petit), #2 curs, #3 tipus de document i unitat, #4 títol
\newcommand{\capcalera}[4]{%
  \gdef\fm@cursnom{#2}%
  \par\noindent
  \begin{tcolorbox}[colback=curs,colframe=curs,boxrule=0pt,arc=2.2mm,width=3.5cm,height=2.15cm,
      valign=center,left=3mm,right=2mm,top=1mm,bottom=1mm,nobeforeafter,box align=center]
    \color{curstext}{\footnotesize #1}\par\vspace{1pt}{\fontsize{19}{21}\selectfont\bfseries #2\par}%
  \end{tcolorbox}\hspace{5mm}%
  \begin{minipage}[c]{\dimexpr\linewidth-3.5cm-5mm-2.9cm\relax}
    {\color{cursfosc}\small\bfseries\MakeUppercase{#3}}\par\vspace{2pt}
    {\color{navy}\fontsize{20}{23}\selectfont\bfseries #4\par}
  \end{minipage}\hfill
  \begin{minipage}[c]{2.6cm}\raggedleft\includegraphics[width=2.6cm]{\fm@logo}\end{minipage}\par
  \vspace{2pt}}

% ---------- barres de secció del color del curs ----------
\renewcommand{\seccio}[2]{\par\needspace{\fmneed\baselineskip}\bigskip
  \begin{tcolorbox}[colback=curs,colframe=curs,coltext=curstext,arc=2mm,boxrule=0pt,left=3mm,right=3mm,top=1.2mm,bottom=1.2mm]
  {\large\bfseries #1\quad #2}
  \end{tcolorbox}\vspace{-2pt}}

% ---------- estrelles de dificultat del color del curs (a Batxillerat no se'n posen) ----------
\@ifundefined{estrelles}{}{%
  \renewcommand{\estrelles}[1]{{\color{curs}\ifcase#1\or\faStar\faStar[regular]\faStar[regular]\or\faStar\faStar\faStar[regular]\or\faStar\faStar\faStar\fi}}}

% ---------- capçalera i peu de les pàgines ----------
% primera pàgina: logos institucionals al peu
\fancypagestyle{primera}{\fancyhf{}\renewcommand{\headrulewidth}{0pt}
  \fancyfoot[L]{\raisebox{-0.32\height}{\includegraphics[height=0.78cm]{logo-gva.png}}\hspace{5mm}%
    \raisebox{-0.32\height}{\includegraphics[height=0.78cm]{logo-vermellar.png}}}
  \fancyfoot[R]{\scriptsize\color{muted}\href{https://creativecommons.org/licenses/by-nc-sa/4.0/deed.ca}{\ccbyncsa}\ \ CC BY-NC-SA 4.0 · fentmates.org\qquad\small\thepage}}
% la resta: el símbol a l'esquerra i el curs a la dreta, amb el títol al centre com abans
\fancyhead[L]{\raisebox{-0.3\height}{\includegraphics[height=0.62cm]{\fm@marca}}}
\fancyhead[R]{\small\bfseries\color{cursfosc}\fm@cursnom}
\renewcommand{\headrule}{\color{curs!45}\hrule height \headrulewidth}
\makeatother

%   \fmcurs{1eso}
% i, al començament del document:
%   \capcalera{Matemàtiques}{1r ESO}{Fitxa d'exercicis · Unitat 2}{Divisibilitat}
% Opcional abans de \input: \def\fmneed{14} (espai mínim abans de cada secció).
% =====================================================================
\makeatletter
\def\fm@logodir{../../logo/pdf}
\providecommand{\fmneed}{10}
\def\fm@cursnom{}

% clau del curs, color i color del text damunt del color
\newcommand{\fm@curs}[3]{\@namedef{fm@color@#1}{#2}\@namedef{fm@text@#1}{#3}}
\fm@curs{1eso}{E6B84A}{navy}
\fm@curs{2eso}{3B82F6}{white}
\fm@curs{3eso}{FF7A50}{white}
\fm@curs{4eso-a}{A78BFA}{white}
\fm@curs{4eso-b}{6D28D9}{white}
\fm@curs{1batx-mi}{9B1C1C}{white}
\fm@curs{1batx-ccss}{EF6B6B}{white}
\fm@curs{2batx-mii}{1E6B47}{white}
\fm@curs{2batx-ccss}{2BB3A3}{white}

\newcommand{\fmcurs}[1]{%
  \@ifundefined{fm@color@#1}{\PackageError{fentmates}{Curs desconegut: #1}{}}{}%
  \edef\fm@tmp{\noexpand\definecolor{curs}{HTML}{\@nameuse{fm@color@#1}}%
               \noexpand\colorlet{curstext}{\@nameuse{fm@text@#1}}}\fm@tmp
  \colorlet{cursfosc}{curs!72!black}%
  \def\fm@logo{\fm@logodir/vertical-#1.pdf}%
  \def\fm@marca{\fm@logodir/marca-#1.pdf}}

% ---------- capçalera de la primera pàgina ----------
% #1 matèria (petit), #2 curs, #3 tipus de document i unitat, #4 títol
\newcommand{\capcalera}[4]{%
  \gdef\fm@cursnom{#2}%
  \par\noindent
  \begin{tcolorbox}[colback=curs,colframe=curs,boxrule=0pt,arc=2.2mm,width=3.5cm,height=2.15cm,
      valign=center,left=3mm,right=2mm,top=1mm,bottom=1mm,nobeforeafter,box align=center]
    \color{curstext}{\footnotesize #1}\par\vspace{1pt}{\fontsize{19}{21}\selectfont\bfseries #2\par}%
  \end{tcolorbox}\hspace{5mm}%
  \begin{minipage}[c]{\dimexpr\linewidth-3.5cm-5mm-2.9cm\relax}
    {\color{cursfosc}\small\bfseries\MakeUppercase{#3}}\par\vspace{2pt}
    {\color{navy}\fontsize{20}{23}\selectfont\bfseries #4\par}
  \end{minipage}\hfill
  \begin{minipage}[c]{2.6cm}\raggedleft\includegraphics[width=2.6cm]{\fm@logo}\end{minipage}\par
  \vspace{2pt}}

% ---------- barres de secció del color del curs ----------
\renewcommand{\seccio}[2]{\par\needspace{\fmneed\baselineskip}\bigskip
  \begin{tcolorbox}[colback=curs,colframe=curs,coltext=curstext,arc=2mm,boxrule=0pt,left=3mm,right=3mm,top=1.2mm,bottom=1.2mm]
  {\large\bfseries #1\quad #2}
  \end{tcolorbox}\vspace{-2pt}}

% ---------- estrelles de dificultat del color del curs (a Batxillerat no se'n posen) ----------
\@ifundefined{estrelles}{}{%
  \renewcommand{\estrelles}[1]{{\color{curs}\ifcase#1\or\faStar\faStar[regular]\faStar[regular]\or\faStar\faStar\faStar[regular]\or\faStar\faStar\faStar\fi}}}

% ---------- capçalera i peu de les pàgines ----------
% primera pàgina: logos institucionals al peu
\fancypagestyle{primera}{\fancyhf{}\renewcommand{\headrulewidth}{0pt}
  \fancyfoot[L]{\raisebox{-0.32\height}{\includegraphics[height=0.78cm]{logo-gva.png}}\hspace{5mm}%
    \raisebox{-0.32\height}{\includegraphics[height=0.78cm]{logo-vermellar.png}}}
  \fancyfoot[R]{\scriptsize\color{muted}\href{https://creativecommons.org/licenses/by-nc-sa/4.0/deed.ca}{\ccbyncsa}\ \ CC BY-NC-SA 4.0 · fentmates.org\qquad\small\thepage}}
% la resta: el símbol a l'esquerra i el curs a la dreta, amb el títol al centre com abans
\fancyhead[L]{\raisebox{-0.3\height}{\includegraphics[height=0.62cm]{\fm@marca}}}
\fancyhead[R]{\small\bfseries\color{cursfosc}\fm@cursnom}
\renewcommand{\headrule}{\color{curs!45}\hrule height \headrulewidth}
\makeatother

%   \fmcurs{1eso}
% i, al començament del document:
%   \capcalera{Matemàtiques}{1r ESO}{Fitxa d'exercicis · Unitat 2}{Divisibilitat}
% Opcional abans de \input: \def\fmneed{14} (espai mínim abans de cada secció).
% =====================================================================
\makeatletter
\def\fm@logodir{../../logo/pdf}
\providecommand{\fmneed}{10}
\def\fm@cursnom{}

% clau del curs, color i color del text damunt del color
\newcommand{\fm@curs}[3]{\@namedef{fm@color@#1}{#2}\@namedef{fm@text@#1}{#3}}
\fm@curs{1eso}{E6B84A}{navy}
\fm@curs{2eso}{3B82F6}{white}
\fm@curs{3eso}{FF7A50}{white}
\fm@curs{4eso-a}{A78BFA}{white}
\fm@curs{4eso-b}{6D28D9}{white}
\fm@curs{1batx-mi}{9B1C1C}{white}
\fm@curs{1batx-ccss}{EF6B6B}{white}
\fm@curs{2batx-mii}{1E6B47}{white}
\fm@curs{2batx-ccss}{2BB3A3}{white}

\newcommand{\fmcurs}[1]{%
  \@ifundefined{fm@color@#1}{\PackageError{fentmates}{Curs desconegut: #1}{}}{}%
  \edef\fm@tmp{\noexpand\definecolor{curs}{HTML}{\@nameuse{fm@color@#1}}%
               \noexpand\colorlet{curstext}{\@nameuse{fm@text@#1}}}\fm@tmp
  \colorlet{cursfosc}{curs!72!black}%
  \def\fm@logo{\fm@logodir/vertical-#1.pdf}%
  \def\fm@marca{\fm@logodir/marca-#1.pdf}}

% ---------- capçalera de la primera pàgina ----------
% #1 matèria (petit), #2 curs, #3 tipus de document i unitat, #4 títol
\newcommand{\capcalera}[4]{%
  \gdef\fm@cursnom{#2}%
  \par\noindent
  \begin{tcolorbox}[colback=curs,colframe=curs,boxrule=0pt,arc=2.2mm,width=3.5cm,height=2.15cm,
      valign=center,left=3mm,right=2mm,top=1mm,bottom=1mm,nobeforeafter,box align=center]
    \color{curstext}{\footnotesize #1}\par\vspace{1pt}{\fontsize{19}{21}\selectfont\bfseries #2\par}%
  \end{tcolorbox}\hspace{5mm}%
  \begin{minipage}[c]{\dimexpr\linewidth-3.5cm-5mm-2.9cm\relax}
    {\color{cursfosc}\small\bfseries\MakeUppercase{#3}}\par\vspace{2pt}
    {\color{navy}\fontsize{20}{23}\selectfont\bfseries #4\par}
  \end{minipage}\hfill
  \begin{minipage}[c]{2.6cm}\raggedleft\includegraphics[width=2.6cm]{\fm@logo}\end{minipage}\par
  \vspace{2pt}}

% ---------- barres de secció del color del curs ----------
\renewcommand{\seccio}[2]{\par\needspace{\fmneed\baselineskip}\bigskip
  \begin{tcolorbox}[colback=curs,colframe=curs,coltext=curstext,arc=2mm,boxrule=0pt,left=3mm,right=3mm,top=1.2mm,bottom=1.2mm]
  {\large\bfseries #1\quad #2}
  \end{tcolorbox}\vspace{-2pt}}

% ---------- estrelles de dificultat del color del curs (a Batxillerat no se'n posen) ----------
\@ifundefined{estrelles}{}{%
  \renewcommand{\estrelles}[1]{{\color{curs}\ifcase#1\or\faStar\faStar[regular]\faStar[regular]\or\faStar\faStar\faStar[regular]\or\faStar\faStar\faStar\fi}}}

% ---------- capçalera i peu de les pàgines ----------
% primera pàgina: logos institucionals al peu
\fancypagestyle{primera}{\fancyhf{}\renewcommand{\headrulewidth}{0pt}
  \fancyfoot[L]{\raisebox{-0.32\height}{\includegraphics[height=0.78cm]{logo-gva.png}}\hspace{5mm}%
    \raisebox{-0.32\height}{\includegraphics[height=0.78cm]{logo-vermellar.png}}}
  \fancyfoot[R]{\scriptsize\color{muted}\href{https://creativecommons.org/licenses/by-nc-sa/4.0/deed.ca}{\ccbyncsa}\ \ CC BY-NC-SA 4.0 · fentmates.org\qquad\small\thepage}}
% la resta: el símbol a l'esquerra i el curs a la dreta, amb el títol al centre com abans
\fancyhead[L]{\raisebox{-0.3\height}{\includegraphics[height=0.62cm]{\fm@marca}}}
\fancyhead[R]{\small\bfseries\color{cursfosc}\fm@cursnom}
\renewcommand{\headrule}{\color{curs!45}\hrule height \headrulewidth}
\makeatother

\fmcurs{4eso-b}

\begin{document}
\thispagestyle{primera}

\capcalera{Matemáticas B}{4\textsuperscript{o} ESO}{Ficha de ejercicios · Unidad 5}{\raggedright Trigonometría}
\vspace{10pt}
Nombre y apellidos: \rule{6.6cm}{0.4pt}\hfill Grupo: \rule{1.6cm}{0.4pt}\hfill Fecha: \rule{2.2cm}{0.4pt}
\vspace{8pt}

\begin{tcolorbox}[colback=soft,colframe=soft,arc=2mm,left=3mm,right=3mm,top=1.5mm,bottom=1.5mm]
\small Ejercicios graduados por niveles de dificultad:\quad \estrelles{1}\ básico \qquad \estrelles{2}\ medio \qquad \estrelles{3}\ avanzado.\\[2pt] Resuélvelos en el cuaderno mostrando todos los pasos y haciendo un dibujo en los problemas. Usa la calculadora en modo \texttt{D} (grados) y da los ángulos en grados, minutos y segundos y las longitudes con dos decimales, salvo que el enunciado pida el resultado exacto. La ficha recoge los ejercicios de las fichas del departamento y añade otros nuevos.
\end{tcolorbox}

% =====================================================
\seccio{1}{Medida de ángulos: grados y radianes}

\exf{1} Expresa en radianes estos ángulos (de manera exacta, en función de $\pi$):\begin{tasks}(6)
  \task $45^\circ$
  \task $30^\circ$
  \task $120^\circ$
  \task $150^\circ$
  \task $540^\circ$
  \task $1080^\circ$
  \task $1800^\circ$
  \task $240^\circ$
  \task $310^\circ$
  \task $70^\circ$
  \task $12^\circ$
  \task $720^\circ$
\end{tasks}

\ex{1} Expresa en grados sexagesimales estos ángulos en radianes:\begin{tasks}(6)
  \task $\pi$
  \task $4\pi$
  \task $7\pi$
  \task $\frac{\pi}{6}$
  \task $\frac{3\pi}{4}$
  \task $\frac{3\pi}{2}$
  \task $\frac{5\pi}{3}$
  \task $\frac{\pi}{2}$
  \task $\frac{4\pi}{3}$
  \task $1{,}2\pi$
  \task $\frac{2\pi}{5}$
  \task $\frac{\pi}{9}$
\end{tasks}

\ex{1} Completa la tabla siguiente:\begin{center}\renewcommand{\arraystretch}{1.6}\begin{tabular}{|l|*{10}{C{1.05cm}|}}\hline\textbf{Grados} & $30^\circ$ &  & $60^\circ$ &  & $135^\circ$ &  & $210^\circ$ &  & $300^\circ$ & \\\hline\textbf{Radianes} &  & $\frac{\pi}{4}$ &  & $\frac{2\pi}{3}$ &  & $\frac{5\pi}{6}$ &  & $\frac{5\pi}{4}$ &  & $\frac{11\pi}{6}$\\\hline\end{tabular}\end{center}

\ex{1} Expresa, tanto en grados sexagesimales como en radianes, estas partes de una circunferencia:\begin{tasks}(4)
  \task Una cuarta parte.
  \task Tres cuartas partes.
  \task Una sexta parte.
  \task Una tercera parte.
\end{tasks}

\ex{2} En el hexágono regular $ABCDEF$ de centro $O$, expresa en radianes los ángulos $\widehat{AOB}$, $\widehat{ABC}$, $\widehat{BAC}$ y $\widehat{ACD}$.\begin{center}\begin{tikzpicture}[scale=1.25]
\foreach \i/\n in {0/D,1/E,2/F,3/A,4/B,5/C} \coordinate (\n) at (60*\i:1.2);
\coordinate (O) at (0,0);
\draw[navy, line width=1.2pt] (A)--(B)--(C)--(D)--(E)--(F)--cycle;
\draw[muted, dashed] (A)--(C) (A)--(D) (O)--(B);
\fill (O) circle (1.3pt) node[above right=-1pt,font=\small]{$O$};
\foreach \n/\p in {A/left,B/below left,C/below right,D/right,E/above right,F/above left} \node[\p,font=\small] at (\n) {$\n$};
\end{tikzpicture}\end{center}

\ex{2} Pasa a grados, minutos y segundos: \ $1$ rad, \ $0{,}75$ rad \ y \ $2{,}5$ rad. Pasa a radianes, de manera exacta: \ $22^\circ30'$, \ $52^\circ30'$ \ y \ $157^\circ30'$.

\ex{3} La aguja de las horas de un reloj da una vuelta en 12 horas. ¿Qué ángulo, en grados y en radianes, recorre en una hora? ¿Y la aguja de los minutos en 20 minutos? ¿Qué ángulo forman las dos agujas a las 4 en punto?


% =====================================================
\seccio{2}{Razones trigonométricas de un ángulo agudo}

\exf{1} Calcula las razones trigonométricas (seno, coseno y tangente) del ángulo agudo $\alpha$ y del otro ángulo agudo, $\beta$, de estos triángulos rectángulos:\begin{enumerate}[label=\alph*),leftmargin=*,itemsep=2pt]\item Catetos de 30 cm y 16 cm; $\alpha$ es opuesto al cateto de 16 cm.\item Hipotenusa de 10 cm y cateto de 6 cm opuesto a $\alpha$.\item Catetos de 24 cm y 10 cm; $\alpha$ es opuesto al cateto de 24 cm.\item Hipotenusa de 25 cm y cateto de 7 cm contiguo a $\alpha$.\item Hipotenusa de 29 cm y cateto de 20 cm contiguo a $\alpha$.\end{enumerate}

\ex{1} Calcula las razones trigonométricas de los ángulos $\hat A$, $\hat C$, $\widehat{ABD}$ y $\widehat{CBD}$.\begin{center}\begin{tikzpicture}[scale=0.27]
\coordinate (A) at (0,0); \coordinate (D) at (9,0); \coordinate (B) at (9,12); \coordinate (C) at (25,0);
\draw[navy, line width=1.1pt] (A)--(B)--(C)--cycle; \draw[muted, dashed] (B)--(D);
\draw (D)++(-0.9,0)--++(0,0.9)--++(0.9,0);
\node[below left] at (A) {$A$}; \node[above] at (B) {$B$}; \node[below right] at (C) {$C$}; \node[below] at (D) {$D$};
\node[above left, font=\small] at (4.5,6) {15 cm}; \node[right, font=\small] at (9,5) {12 cm}; \node[below, font=\small] at (17,0) {16 cm};
\end{tikzpicture}\end{center}

\ex{1} Completa la tabla con la calculadora (con cuatro decimales):\begin{center}\renewcommand{\arraystretch}{1.45}\begin{tabular}{|c|*{4}{C{2.3cm}|}}\hline $\alpha$ & $15^\circ$ & $55^\circ20'$ & $72^\circ25'40''$ & $85{,}5^\circ$\\\hline
$\sin\alpha$ & & & & \\\hline $\cos\alpha$ & & & & \\\hline $\tg\alpha$ & & & & \\\hline\end{tabular}\end{center}

\ex{1} Halla, con la calculadora, el ángulo agudo $\alpha$ en cada caso:\begin{tasks}(5)
  \task $\sin\alpha=\frac{\sqrt3}{2}$
  \task $\cos\alpha=\frac13$
  \task $\tg\alpha=2{,}5$
  \task $\sin\alpha=0{,}25$
  \task $\cos\alpha=0{,}8$
\end{tasks}

\ex{2} Comprueba, para $\alpha=60^\circ$, a partir de los hexágonos regulares de lado 4 cm y de lado 6 cm (formados por seis triángulos equiláteros), que las razones trigonométricas no varían según la longitud del lado.

\ex{2} Calcula los lados que faltan en cada triángulo rectángulo y, si hace falta, las otras razones del ángulo:\begin{enumerate}[label=\alph*),leftmargin=*,itemsep=2pt]\item El cateto contiguo a un ángulo $\alpha$ mide 5 cm y $\tg\alpha=\frac65$.\item El seno de un ángulo vale $\frac{12}{13}$ y el cateto opuesto mide 24 cm. Calcula también el coseno y la tangente.\item El cateto opuesto a un ángulo mide 15 cm y la tangente de ese ángulo vale $\frac34$. Calcula también el seno y el coseno.\item La hipotenusa mide 2 cm y la tangente de uno de los ángulos vale 1.\end{enumerate}

\ex{3} La hipotenusa de un triángulo rectángulo mide 10 cm y la tangente de uno de los ángulos agudos vale $\frac{7}{12}$. Calcula los catetos y el seno y el coseno de ese ángulo.


% =====================================================
\seccio{3}{Razones trigonométricas de 30°, 45° y 60°}

\exf{1} Completa la tabla siguiente sin calculadora, de manera exacta:\begin{center}\renewcommand{\arraystretch}{1.6}\begin{tabular}{|l|*{5}{C{1.6cm}|}}\hline
& $0^\circ$ & $30^\circ$ & $45^\circ$ & $60^\circ$ & $90^\circ$\\\hline
$\sin$ & & & & & \\\hline $\cos$ & & & & & \\\hline $\tg$ & & & & & \\\hline\end{tabular}\end{center}

\ex{1} En cada triángulo rectángulo se conocen un ángulo agudo y un lado. Calcula los otros dos lados de manera exacta:\begin{enumerate}[label=\alph*),leftmargin=*,itemsep=2pt]\item Un ángulo de $60^\circ$ y el cateto opuesto, de 5 cm.\item Un ángulo de $30^\circ$ y el cateto opuesto, de 2 cm.\item Un ángulo de $45^\circ$ y el cateto contiguo, de 7 cm.\item Un ángulo de $60^\circ$ y el cateto contiguo, de 6 cm.\item Un ángulo de $45^\circ$ y el cateto opuesto, de 12 cm.\item Un ángulo de $30^\circ$ y la hipotenusa, de 30 cm.\end{enumerate}

\ex{1} Calcula sin calculadora:\begin{tasks}(3)
  \task $\sin30^\circ+\cos60^\circ$
  \task $\tg45^\circ-2\sin^245^\circ$
  \task $\sin60^\circ\cos30^\circ-\sin30^\circ\cos60^\circ$
  \task $\tg60^\circ\cdot\tg30^\circ$
  \task $(\sin45^\circ+\cos45^\circ)^2$
  \task $\cos^230^\circ-\sin^230^\circ$
\end{tasks}

\ex{2} Sin utilizar el teorema de Pitágoras, calcula la altura de un triángulo equilátero de 8 cm de lado y la diagonal de un cuadrado de 5 cm de lado.

\ex{2} Un hexágono regular se puede dividir en 6 triángulos equiláteros. Si el lado del hexágono mide 4 cm, calcula su apotema y su área de manera exacta.

\ex{3} Dos antenas de radio, de 100 m y 75 m, están sujetas al suelo como se indica en la figura. Calcula la longitud de cada tramo de cable y la distancia $\overline{AE}$. Da los resultados de manera exacta.\begin{center}\begin{tikzpicture}[scale=0.62]
\coordinate (A) at (0,0); \coordinate (P) at (2.3,0); \coordinate (B) at (2.3,4); \coordinate (C) at (9.2,0); \coordinate (Q) at (12.2,0); \coordinate (D) at (12.2,3); \coordinate (E) at (13.9,0);
\draw[muted] (-0.6,0)--(14.6,0);
\draw[accent, line width=1.6pt] (P)--(B) (Q)--(D);
\draw[navy, dashed] (A)--(B)--(C)--(D)--(E);
\foreach \n/\p in {A/below,P/below,C/below,Q/below,E/below,B/above,D/above} \node[\p,font=\small] at (\n) {$\n$};
\node[left,font=\small] at (2.3,2) {100 m}; \node[left,font=\small] at (12.2,1.4) {75 m};
\draw (0.8,0) arc (0:60:0.8); \node[font=\scriptsize] at (1.15,0.4) {$60^\circ$};
\draw (8.2,0) arc (180:150:1); \node[font=\scriptsize] at (7.7,0.3) {$30^\circ$};
\draw (12.2,2.2) arc (270:225:0.8); \node[font=\scriptsize] at (11.7,1.9) {$45^\circ$};
\draw (12.2,2.2) arc (270:300:0.8); \node[font=\scriptsize] at (12.75,1.85) {$30^\circ$};
\end{tikzpicture}\end{center}


% =====================================================
\seccio{4}{Relaciones fundamentales}

\exf{1} Sabiendo que el ángulo $\alpha$ es agudo, calcula, en cada caso, sus otras razones trigonométricas de manera exacta:\begin{tasks}(5)
  \task $\sin\alpha=\frac{7}{25}$
  \task $\sin\alpha=\frac{12}{37}$
  \task $\sin\alpha=\frac{11}{61}$
  \task $\cos\alpha=\frac{20}{29}$
  \task $\cos\alpha=\frac{12}{13}$
  \task $\sin\alpha=\frac{3}{8}$
  \task $\cos\alpha=\frac{1}{3}$
  \task $\sin\alpha=\frac{15}{34}$
  \task $\cos\alpha=\frac{2}{5}$
  \task $\sin\alpha=\frac{4}{7}$
\end{tasks}\recorda{$\sin^2\alpha+\cos^2\alpha=1$, \quad $\tg\alpha=\dfrac{\sin\alpha}{\cos\alpha}$, \quad $\tg^2\alpha+1=\dfrac{1}{\cos^2\alpha}$.}

\ex{2} Si el ángulo $\alpha$ es agudo, calcula, en cada caso, el seno y el coseno del ángulo de manera exacta:\begin{tasks}(7)
  \task $\tg\alpha=\frac{5}{12}$
  \task $\tg\alpha=\frac{3}{4}$
  \task $\tg\alpha=\frac{5}{3}$
  \task $\tg\alpha=\frac{5}{13}$
  \task $\tg\alpha=\frac{5}{8}$
  \task $\tg\alpha=3$
  \task $\tg\alpha=1$
\end{tasks}

\ex{1} Sabiendo que el ángulo $\alpha$ es agudo, calcula las otras dos razones (con cuatro decimales):\begin{tasks}(3)
  \task $\cos\alpha=0,1357$
  \task $\cos\alpha=0,9531$
  \task $\cos\alpha=0,2864$
  \task $\sin\alpha=0,1827$
  \task $\sin\alpha=0,9542$
  \task $\sin\alpha=0,4531$
\end{tasks}

\ex{2} Comprueba si hay algún ángulo que tenga estas razones trigonométricas:\begin{tasks}(2)
  \task $\sin\alpha=\frac{7}{25}$ \ y \ $\cos\alpha=\frac{11}{25}$
  \task $\sin\alpha=\frac{7}{13}$ \ y \ $\cos\alpha=\frac{12}{13}$
  \task $\sin\alpha=\frac{\sqrt{5}}{5}$ \ y \ $\cos\alpha=\frac{2 \sqrt{5}}{5}$
  \task $\sin\alpha=\frac{3}{5}$ \ y \ $\cos\alpha=\frac{4}{5}$
  \task $\sin\alpha=\frac{7}{11}$ \ y \ $\tg\alpha=\frac{6}{11}$
  \task $\cos\alpha=\frac{5}{13}$ \ y \ $\tg\alpha=1$
  \task $\cos\alpha=\frac{\sqrt{3}}{2}$ \ y \ $\tg\alpha=\frac{\sqrt{3}}{3}$
  \task $\cos\alpha=\frac{3}{5}$ \ y \ $\tg\alpha=\frac{7}{5}$
\end{tasks}

\ex{2} Comprueba si estos pares de razones trigonométricas (redondeadas a cuatro decimales) pueden ser del mismo ángulo:\begin{tasks}(2)
  \task $\cos\alpha=0,1736$ \ y \ $\tg\alpha=0,4663$
  \task $\cos\alpha=0,9397$ \ y \ $\tg\alpha=0,3640$
  \task $\sin\alpha=0,1485$ \ y \ $\cos\alpha=0,2588$
  \task $\cos\alpha=0,7313$ \ y \ $\tg\alpha=0,9325$
  \task $\sin\alpha=0,2754$ \ y \ $\cos\alpha=0,6691$
\end{tasks}\recorda{con decimales redondeados, los dos miembros de la relación pueden diferir un poco en la última cifra.}

\ex{3} Utiliza las relaciones fundamentales para demostrar estas igualdades:\begin{tasks}(1)
  \task $\dfrac{1}{\tg^2\alpha}+1=\dfrac{1}{\sin^2\alpha}$
  \task $(1+\tg^2\alpha)\cdot\cos^2\alpha=1$
  \task $\sin\alpha\cdot\cos\alpha\cdot\tg\alpha+\cos^2\alpha=1$
\end{tasks}


% =====================================================
\seccio{5}{Resolución de triángulos rectángulos}

\exf{1} Resuelve los triángulos rectángulos siguientes. En todos, $\hat A=90^\circ$; por tanto, $a$ es la hipotenusa.\begin{tasks}(2)
  \task $b=4$ m, \ $\hat B=30^\circ$
  \task $a=7$ m, \ $b=5$ m
  \task $a=10$ m, \ $\hat B=40^\circ$
  \task $c=5$ cm, \ $\sin\hat C=0{,}25$
\end{tasks}

\ex{1} Calcula la altura $h$, trazada desde el vértice $B$, de cada triángulo $ABC$:\begin{tasks}(2)
  \task $\overline{AB}=18$ cm y $\hat A=62^\circ$.
  \task $\overline{BC}=28$ cm y $\hat C=35^\circ$.
\end{tasks}

\ex{1} Una escalera de 3 m está apoyada en una pared. ¿Qué ángulo forma la escalera con el suelo si la base está a 1,2 m de la pared?

\ex{1} Una torre de 20 m de altura proyecta una sombra de 25 m de longitud. Calcula la inclinación de los rayos solares en ese momento.

\ex{2} En un triángulo isósceles, el lado desigual mide 18 m y la altura, 10 m. ¿Cuánto miden sus ángulos?

\ex{2} En una ruta de montaña, una señal indica una altitud de 785 m. Tres kilómetros más adelante (en horizontal), la altitud es de 1\,265 m. Halla la pendiente media de la ruta y el ángulo que forma con la horizontal.

\ex{2} Un avión que está a punto de aterrizar desciende 2\,160 m al recorrer una distancia de 18,41 km. ¿Cuál es el ángulo $\alpha$ con el que desciende?

\ex{2} Los brazos de un compás, que miden 12 cm, forman un ángulo de $50^\circ$. ¿Cuál es el radio de la circunferencia que se puede trazar con esa abertura?

\ex{2} Dos amigos salen del punto $A$. Uno se dirige al punto $B$ y el otro al punto $C$, siguiendo trayectorias rectilíneas que forman entre sí un ángulo de $37^\circ$. Si los puntos $B$ y $C$ están separados 120 m y la dirección $BC$ es perpendicular a la dirección $AC$, ¿cuántos metros recorre uno de los amigos más que el otro?


% =====================================================
\seccio{6}{Problemas de trigonometría}

\exf{2} Desde cierto lugar se ve el punto más alto de una torre con un ángulo de elevación de $35^\circ$. Si retrocedemos 200 m, se ve la misma torre con un ángulo de $20^\circ$. Calcula la altura de la torre.

\ex{2} Para calcular la altura de un edificio se miden los ángulos de elevación desde dos puntos alineados con el pie del edificio y situados a una distancia de 100 m el uno del otro. ¿Cuál es la altura del edificio si los ángulos son $33^\circ$ y $46^\circ$?

\ex{2} Para medir la altura de una montaña se miden los ángulos de elevación desde dos puntos alineados con la cima, situados a una distancia de 480 m el uno del otro y a una altitud de 1\,200 m sobre el nivel del mar. ¿Cuál es la altura de la montaña si los ángulos son $45^\circ$ y $76^\circ$?

\ex{2} Para hallar la altura a la que se encuentra un globo, Rosa se coloca en un punto $B$, y yo en un punto $A$, a 5 metros de ella, de manera que $A$, $B$ y el pie $C$ de la vertical del globo quedan alineados ($B$ entre $A$ y $C$). Si desde $A$ veo el globo con un ángulo de $40^\circ$ y Rosa lo ve con un ángulo de $50^\circ$, ¿a qué altura se encuentra el globo?

\ex{2} Un avión vuela entre dos ciudades, $A$ y $B$, que están a 75 km la una de la otra. Las visuales desde $A$ y desde $B$ hacia el avión forman con la horizontal ángulos de $36^\circ$ y $12^\circ$, respectivamente. Calcula la altura a la que vuela el avión y las distancias a las que se encuentra de $A$ y de $B$ (el avión y las ciudades están en el mismo plano vertical).

\ex{2} Queremos colgar una lámpara del techo de una habitación con una cuerda atada en dos puntos del techo separados 140 cm, de manera que los ángulos entre la cuerda y el techo sean de $60^\circ$ y $40^\circ$. ¿Cuál será la longitud total de la cuerda? ¿A qué distancia del techo quedará la lámpara?\begin{center}\begin{tikzpicture}[scale=0.032]
\draw[navy, line width=1.2pt] (-10,0)--(150,0); \foreach \x in {-6,4,...,146} \draw[navy] (\x,0)--++(4,4);
\coordinate (P) at (0,0); \coordinate (Q) at (140,0); \coordinate (L) at (45.69,-79.14);
\draw[muted, line width=0.9pt] (P)--(L)--(Q); \draw[muted, dashed] (L)--(45.69,0);
\fill[accent] (L) circle (2.6);
\draw (12,0) arc (0:-60:12); \node[font=\scriptsize] at (19,-9) {$60^\circ$};
\draw (128,0) arc (180:220:12); \node[font=\scriptsize] at (117,-7) {$40^\circ$};
\draw[<->, muted] (0,14)--(140,14); \node[above, font=\small] at (70,14) {140 cm};
\end{tikzpicture}\end{center}

\ex{2} Una antena de radio está sujeta al suelo por dos cables de acero, como indica la figura. Calcula, de manera exacta, la altura de la antena, la longitud de los cables y el ángulo $\widehat{ABC}$.\begin{center}\begin{tikzpicture}[scale=0.055]
\coordinate (A) at (0,0); \coordinate (C) at (126,0); \coordinate (H) at (31.5,0); \coordinate (B) at (31.5,54.56);
\draw[muted] (-6,0)--(132,0); \draw[navy] (A)--(B)--(C); \draw[accent, line width=1.6pt] (H)--(B);
\node[left] at (A) {$A$}; \node[right] at (C) {$C$}; \node[above] at (B) {$B$};
\draw (9,0) arc (0:60:9); \node[font=\scriptsize] at (13,6) {$60^\circ$};
\draw (114,0) arc (180:150:12); \node[font=\scriptsize] at (104,3.5) {$30^\circ$};
\draw[<->, muted] (0,-6)--(126,-6); \node[below, font=\small] at (63,-6) {126 m};
\end{tikzpicture}\end{center}

\ex{2} Una escalera de bomberos de 10 m de longitud se ha fijado en un punto de una calle. Si se apoya sobre una de las fachadas, forma con el suelo un ángulo de $45^\circ$, y si se apoya sobre la otra fachada, el ángulo es de $30^\circ$. ¿Cuál es la anchura de la calle?

\ex{2} Una escalera para acceder a un túnel tiene la forma y las dimensiones de la figura. Calcula la profundidad a la que está el punto $B$ respecto del punto $A$.\begin{center}\begin{tikzpicture}[scale=0.11]
\coordinate (A) at (0,0); \coordinate (P) at ($(A)+(-30:25)$); \coordinate (Q) at ($(P)+(10,0)$); \coordinate (B) at ($(Q)+(-50:30)$);
\draw[navy, line width=1.6pt, line join=round] (A)--(P)--(Q)--(B);
\draw[muted, dashed] (P)--++(-7,0) (B)--++(-8,0);
\node[above left] at (A) {$A$}; \node[below right] at (B) {$B$};
\node[above right, font=\small] at ($(A)!0.5!(P)$) {25 m}; \node[above, font=\small] at ($(P)!0.5!(Q)$) {10 m}; \node[right, font=\small] at ($(Q)!0.5!(B)+(0.6,0)$) {30 m};
\draw ($(P)+(-4,0)$) arc (180:150:4); \node[font=\scriptsize] at ($(P)+(-6.5,1.4)$) {$30^\circ$};
\draw ($(B)+(-4,0)$) arc (180:130:4); \node[font=\scriptsize] at ($(B)+(-6.5,2.4)$) {$50^\circ$};
\end{tikzpicture}\end{center}

\ex{2} Polígonos regulares:\begin{enumerate}[label=\alph*),leftmargin=*,itemsep=2pt]\item ¿Cuánto mide la apotema de un pentágono regular de 10 cm de lado?\item Calcula el área de un pentágono regular inscrito en una circunferencia de 6 cm de radio.\end{enumerate}

\ex{2} Triángulos isósceles:\begin{enumerate}[label=\alph*),leftmargin=*,itemsep=2pt]\item La base de un triángulo isósceles mide 5 cm y el ángulo opuesto a esa base es de $55^\circ$. Calcula la altura y el área del triángulo.\item El lado desigual de un triángulo isósceles mide 34 cm y el ángulo opuesto, $50^\circ$. Calcula el perímetro del triángulo y su área.\end{enumerate}

\ex{2} Rombos:\begin{enumerate}[label=\alph*),leftmargin=*,itemsep=2pt]\item Calcula el área de un rombo de 12 cm de lado sabiendo que uno de sus ángulos mide $43^\circ$.\item Halla los ángulos de un rombo cuyas diagonales miden 12 cm y 8 cm. Calcula también su perímetro y su área.\end{enumerate}

\ex{3} Calcula la longitud de la diagonal menor y el área de un paralelogramo de lados 10 cm y 15 cm, sabiendo que uno de sus ángulos mide $35^\circ$.

\ex{3} Para construir un túnel de $A$ hasta $B$ se localiza una roca $C$ visible desde los dos puntos. Se miden las distancias $\overline{AC}=245$ m y $\overline{BC}=658$ m, y el ángulo $\widehat{ACB}=57^\circ$. ¿Cuál será la longitud del túnel?

\ex{3} Calcula el ángulo de tiro $\alpha$ del jugador que está en el punto $A$, en el centro de la línea de banda, si la portería $BC$ mide 5 m y está centrada en la línea de fondo.\begin{center}\begin{tikzpicture}[scale=0.055]
\fill[teal!15] (0,0) rectangle (100,60); \draw[navy, line width=1pt] (0,0) rectangle (100,60);
\draw[navy!50] (50,0)--(50,60) (50,30) circle (9);
\coordinate (A) at (50,0); \coordinate (B) at (100,27.5); \coordinate (C) at (100,32.5);
\draw[accent, line width=0.9pt] (A)--(B) (A)--(C); \fill (A) circle (0.9) (B) circle (0.9) (C) circle (0.9);
\node[below] at (A) {$A$}; \node[right] at (B) {$B$}; \node[right] at (C) {$C$};
\node[font=\scriptsize] at (84,23) {$\alpha$};
\draw[<->, muted] (0,64)--(100,64); \node[above, font=\small] at (50,64) {100 m};
\draw[<->, muted] (-4,0)--(-4,60); \node[left, font=\small] at (-4,30) {60 m};
\end{tikzpicture}\end{center}


% =====================================================
\seccio{7}{Razones trigonométricas de un ángulo cualquiera}

\exf{1} Utiliza la calculadora para calcular las razones trigonométricas de estos ángulos (con cuatro decimales) e indica en qué cuadrante está cada uno:\begin{center}\renewcommand{\arraystretch}{1.35}\begin{tabular}{|c|C{2.2cm}|C{2.2cm}|C{2.2cm}|C{2cm}|}\hline $\alpha$ & $\sin\alpha$ & $\cos\alpha$ & $\tg\alpha$ & \textbf{Cuadrante}\\\hline $120^\circ$ & & & & \\\hline $220^\circ$ & & & & \\\hline $135^\circ$ & & & & \\\hline $300^\circ$ & & & & \\\hline $240^\circ$ & & & & \\\hline $75^\circ$ & & & & \\\hline \end{tabular}\end{center}

\ex{1} Responde:\begin{tasks}(1)
  \task Si $\sin\alpha=-0{,}8$, ¿en qué cuadrante puede estar el ángulo $\alpha$?
  \task Si $\cos\beta=-0{,}3$, ¿en qué cuadrante puede estar el ángulo $\beta$?
  \task Si $\tg\gamma=-1{,}5$, ¿en qué cuadrante puede estar el ángulo $\gamma$?
\end{tasks}

\ex{1} Calcula sin calculadora:\begin{tasks}(2)
  \task $3\sin270^\circ-2\cos180^\circ$
  \task $\cos0^\circ+\sin180^\circ-\tg180^\circ$
  \task $\sin90^\circ\cdot\cos360^\circ$
  \task $2\cos90^\circ-\sin270^\circ$
\end{tasks}

\ex{2} Calcula las otras razones trigonométricas del ángulo en cada caso:\begin{tasks}(2)
  \task $\cos\alpha=0{,}7$ \ y \ $\alpha$ en el primer cuadrante.
  \task $\tg\beta=-3$ \ y \ $\beta$ en el segundo cuadrante.
  \task $\sin\gamma=-\frac35$ \ y \ $\gamma$ en el tercer cuadrante.
  \task $\cos\delta=\frac{5}{13}$ \ y \ $\delta$ en el cuarto cuadrante.
\end{tasks}

\ex{2} Calcula las razones trigonométricas de $\alpha$ sabiendo que $\sin\alpha=0{,}34$. ¿Cuántos casos hay?

\ex{2} Halla los dos ángulos de $[0^\circ,360^\circ)$ que cumplen cada condición:\begin{tasks}(4)
  \task $\sin x=-\frac12$
  \task $\tg x=-1$
  \task $\cos x=0{,}8$
  \task $\sin x=0{,}9$
\end{tasks}\recorda{la calculadora da un solo ángulo. El otro es $180^\circ-x$ (seno), $360^\circ-x$ (coseno) o $x+180^\circ$ (tangente).}

\ex{2} Un ángulo $\beta$ tal que $0^\circ<\beta<360^\circ$ verifica $\sin\beta=-\sin30^\circ$ y $\cos\beta=\cos30^\circ$. ¿A qué cuadrante pertenece? ¿Cuánto mide?

\ex{1} ¿A qué ángulo entre $0^\circ$ y $360^\circ$ corresponde cada uno de estos ángulos?\begin{tasks}(4)
  \task $1750^\circ$
  \task $1920^\circ$
  \task $-800^\circ$
  \task $3825^\circ$
\end{tasks}


% =====================================================
\seccio{8}{Reducción al primer cuadrante}

\exf{1} Utiliza la relación entre los ángulos para calcular, de manera exacta:\begin{tasks}(1)
  \task Las razones de $135^\circ$, sabiendo las de $45^\circ$.
  \task Las razones de $240^\circ$, sabiendo las de $60^\circ$.
  \task Las razones de $-30^\circ$, sabiendo las de $30^\circ$.
\end{tasks}

\ex{2} Calcula las razones trigonométricas (seno, coseno y tangente) de los ángulos siguientes, relacionándolas con las de un ángulo del primer cuadrante. Si el ángulo del primer cuadrante es $30^\circ$, $45^\circ$ o $60^\circ$, da su valor exacto.\begin{tasks}(6)
  \task $155^\circ$
  \task $240^\circ$
  \task $390^\circ$
  \task $225^\circ$
  \task $1920^\circ$
  \task $-30^\circ$
\end{tasks}

\ex{2} Expresa las razones trigonométricas de los ángulos siguientes en función de las de un ángulo del primer cuadrante. ¡Cuidado con los signos!\begin{tasks}(5)
  \task $-40^\circ$
  \task $385^\circ$
  \task $-185^\circ$
  \task $835^\circ$
  \task $934^\circ$
\end{tasks}

\ex{2} Sabiendo que $\sin35^\circ\approx0{,}57$, $\cos35^\circ\approx0{,}82$ y $\tg35^\circ\approx0{,}70$, calcula sin calculadora las razones trigonométricas de:\begin{tasks}(6)
  \task $145^\circ$
  \task $215^\circ$
  \task $325^\circ$
  \task $55^\circ$
  \task $-35^\circ$
  \task $395^\circ$
\end{tasks}

\ex{3} Expresa las razones trigonométricas de los ángulos siguientes en función de las de un ángulo $\alpha$ tal que $0^\circ\le\alpha\le45^\circ$:\begin{tasks}(4)
  \task $120^\circ$
  \task $430^\circ$
  \task $-800^\circ$
  \task $180^\circ$
\end{tasks}\recorda{si el ángulo del primer cuadrante es mayor de $45^\circ$, pasa al complementario: $\sin(90^\circ-\alpha)=\cos\alpha$, $\cos(90^\circ-\alpha)=\sin\alpha$, $\tg(90^\circ-\alpha)=\dfrac{1}{\tg\alpha}$.}

\ex{3} Simplifica: \ $\sin(180^\circ-\alpha)+\sin(180^\circ+\alpha)+\cos(360^\circ-\alpha)+\cos(180^\circ-\alpha)$.

\vfill
\begin{center}\color{muted}\small\faCheckCircle[regular]\ \ ¡Ánimo! Comprueba tus resultados con las herramientas interactivas de los apuntes de la unidad.\end{center}

\end{document}
